\subsection{From model-based dynamic programming to learning}

All the algorithmic steps introduced for MDPs relied on model information, in particular on the
transition probabilities
\[
P^u_{x,x'}
=
\Prob[x_{k+1}=x'\mid x_k=x,\ u_k=u].
\]
For instance, in policy iteration, the policy evaluation step required the closed-loop transition
matrix
\[
P^\pi_{x,x'}
=
\sum_{u\in\cU}\pi(x,u)\textcolor{red}{P^u_{x,x'}},
\]
and the value was computed as
\[
V^\pi
=
(I-\gamma P^\pi)^{-1}R^\pi.
\]
Here we have switched notation from costs to rewards. Therefore \(R\) and \(r\) denote rewards,
and the objective is now to maximize value functions instead of minimizing cost functions.

The policy improvement step also required the transition probabilities. Indeed, once \(V^\pi\)
is known, a greedy improvement is obtained by solving
\[
\pi_{\mathrm{new}}(x,\cdot)
\in
\argmax_{\nu(x,\cdot)\in\Delta(\cU)}
\sum_{u\in\cU}\nu(x,u)
\left[
R_x^u
+
\gamma\sum_{x'\in\cX}\textcolor{red}{P^u_{x,x'}}V^\pi(x')
\right].
\]
Equivalently, when a deterministic maximizer is selected,
\[
\mu_{\mathrm{new}}(x)
\in
\argmax_{u\in\cU}
\left[
R_x^u
+
\gamma\sum_{x'\in\cX}\textcolor{red}{P^u_{x,x'}}V^\pi(x')
\right].
\]

The same issue appears in value iteration. The Bellman optimality update is
\[
V(x)
\leftarrow
\max_{\pi(x,\cdot)}
\sum_{u\in\cU}\pi(x,u)
\left[
R_x^u
+
\gamma\sum_{x'\in\cX}\textcolor{red}{P^u_{x,x'}}V(x')
\right],
\]
or, equivalently,
\[
V(x)
\leftarrow
\max_{u\in\cU}
\left[
R_x^u
+
\gamma\sum_{x'\in\cX}\textcolor{red}{P^u_{x,x'}}V(x')
\right].
\]
Thus, both policy iteration and value iteration are model-based algorithms: they require
knowing how likely every next state \(x'\) is after applying every action \(u\) in every state
\(x\).

The central question in learning is therefore the following: can we find a good, or even optimal,
policy without explicitly knowing the transition probabilities?

There are two different settings. In the first setting, we have access to collected data, typically
in the form of repeated episodes,
\[
(x_0,u_0,r_0),\ (x_1,u_1,r_1),\ \dots,\ (x_T,u_T,r_T).
\]
These episodes may come from repeated simulations of the control task, or from repeated
experiments in a controlled environment. In the second setting, learning happens online during
the control task itself, with no prior training. This is much harder, because the controller must
learn and control at the same time. This is the classical dual-control difficulty: the controller
must collect information while also trying to achieve good performance.

In this chapter we focus on the first viewpoint: learning from episodes. The goal is to replace
explicit model-based policy evaluation with empirical estimates obtained from data.

\subsection{The Q-function}

A useful reformulation is obtained by introducing the action-value function, or Q-function. For
a fixed policy \(\pi\), the Q-function is defined as the expected discounted future reward obtained
by:
\begin{itemize}
    \item starting in state \(x\);
    \item taking action \(u\) immediately;
    \item applying policy \(\pi\) at all subsequent times.
\end{itemize}
Therefore,
\[
Q^\pi(x,u)
=
R_x^u
+
\gamma\sum_{x'\in\cX}P^u_{x,x'}V^\pi(x').
\]
The difference between \(V^\pi\) and \(Q^\pi\) is the following. The value function \(V^\pi(x)\)
is defined only over states: it measures the expected return when one starts from \(x\) and
then follows policy \(\pi\). The Q-function \(Q^\pi(x,u)\), instead, is defined over state-action
pairs: it measures the expected return when the first action is fixed to be \(u\), and the policy
\(\pi\) is followed only afterward.

Thus, if the MDP has \(N\) states and \(M\) actions, then
\[
V^\pi
\quad\text{has }N\text{ entries},
\qquad
Q^\pi
\quad\text{has }NM\text{ entries}.
\]
The two objects are directly related. Since under policy \(\pi\), action \(u\) is selected with
probability \(\pi(x,u)\), we have
\[
V^\pi(x)
=
\sum_{u\in\cU}\pi(x,u)Q^\pi(x,u).
\]
In words, the value of a state is the average of the Q-values of all possible actions, weighted
by the probabilities with which the policy selects those actions.

\subsection{Bellman equations in Q-form}

The standard Bellman equation for the value function is
\[
V^\pi(x)
=
\sum_{u\in\cU}\pi(x,u)
\left[
R_x^u
+
\gamma\sum_{x'\in\cX}P^u_{x,x'}V^\pi(x')
\right].
\]
Using the relation
\[
V^\pi(x')
=
\sum_{u'\in\cU}\pi(x',u')Q^\pi(x',u'),
\]
we can write the Bellman equation directly in terms of \(Q^\pi\):
\[
Q^\pi(x,u)
=
R_x^u
+
\gamma\sum_{x'\in\cX}P^u_{x,x'}
\sum_{u'\in\cU}\pi(x',u')Q^\pi(x',u').
\]
This is again a consistency equation. The Q-value of \((x,u)\) is equal to the immediate
expected reward \(R_x^u\), plus the discounted expected Q-value of the next state-action pair.
The next state \(x'\) is generated by the transition probabilities, while the next action \(u'\)
is generated by the policy \(\pi\).

The model information required is essentially the same as before: one still needs
\(P^u_{x,x'}\) and \(R_x^u\). However, the advantage of the Q-function is that policy
improvement becomes much simpler. If \(Q^\pi\) is known, improving the policy does not require
one-step model predictions anymore. It only requires comparing the values \(Q^\pi(x,u)\) over
the actions available at the current state.

\subsection{Bellman optimality in Q-form}

In reward maximization, the optimal value function satisfies
\[
V^\star(x)
=
\max_{u\in\cU}
\left[
R_x^u
+
\gamma\sum_{x'\in\cX}P^u_{x,x'}V^\star(x')
\right].
\]
The corresponding optimal Q-function satisfies
\[
Q^\star(x,u)
=
R_x^u
+
\gamma\sum_{x'\in\cX}P^u_{x,x'}
\left(
\max_{u'\in\cU}Q^\star(x',u')
\right).
\]
Indeed, after taking action \(u\) in state \(x\), the next state \(x'\) is random. Once the system
has reached \(x'\), optimal behavior means choosing an action \(u'\) that maximizes the
optimal Q-value at \(x'\). Therefore,
\[
V^\star(x')
=
\max_{u'\in\cU}Q^\star(x',u').
\]
Equivalently,
\[
Q^\star(x,u)
=
R_x^u
+
\gamma\sum_{x'\in\cX}P^u_{x,x'}V^\star(x').
\]

The important structural change is that the maximization over the next action has moved
inside the Q-form Bellman equation. The expectation is still with respect to the next state, but
the optimal action at the next state is chosen by maximizing \(Q^\star(x',u')\). Once
\(Q^\star\) is known, the optimal policy is immediately obtained as
\[
\mu^\star(x)
\in
\argmax_{u\in\cU}Q^\star(x,u).
\]
Thus the Q-function implicitly defines the optimal strategy.

\subsection{Policy iteration with Q-functions}

Let us now rewrite policy iteration using the Q-function. In the value-function formulation,
policy iteration had two steps:
\begin{enumerate}
    \item compute \(V^\pi\) for the current policy \(\pi\);
    \item improve the policy greedily by using the model and the value \(V^\pi\).
\end{enumerate}
The improvement step required the expression
\[
R_x^u+\gamma\sum_{x'\in\cX}P^u_{x,x'}V^\pi(x').
\]
Therefore, model information entered both the policy evaluation step and the policy improvement
step.

In the Q-function formulation, the first step becomes the computation of \(Q^\pi\), which is
defined by the Bellman equation
\[
Q^\pi(x,u)
=
R_x^u
+
\gamma\sum_{x'\in\cX}P^u_{x,x'}
\sum_{u'\in\cU}\pi(x',u')Q^\pi(x',u').
\]
Once \(Q^\pi\) is known, the greedy policy update is simply
\[
\pi_{\mathrm{new}}(x,\cdot)
\in
\argmax_{\nu(x,\cdot)\in\Delta(\cU)}
\sum_{u\in\cU}\nu(x,u)Q^\pi(x,u).
\]
Equivalently, for a deterministic policy,
\[
\mu_{\mathrm{new}}(x)
\in
\argmax_{u\in\cU}Q^\pi(x,u).
\]

This is the key advantage of Q-functions in a model-free setting. The policy evaluation step is
still difficult, because computing \(Q^\pi\) exactly requires model information and involves
\(NM\) unknowns. However, the policy improvement step becomes trivial and requires no model
information: for each state, one simply chooses the action with the largest estimated Q-value.

Therefore, if we can learn or estimate \(Q^\pi\) from data, we can improve the policy without
ever explicitly estimating the transition probabilities.

\subsection{Experimental policy evaluation}

The idea of Monte Carlo policy evaluation is to delegate the policy evaluation step to
experiments or simulations. Suppose that the system is controlled by a fixed policy \(\pi\), and
that we collect one long episode
\[
(x_0,u_0,r_0),\ (x_1,u_1,r_1),\ \dots,\ (x_T,u_T,r_T).
\]
For each time \(k\), define the empirical discounted return
\[
g_k
=
\sum_{i=k}^{T}\gamma^{i-k}r_i.
\]
This is the reward accumulated from time \(k\) until the end of the episode, with discounting.
It is a sample realization of the return obtained after observing the state-action pair
\((x_k,u_k)\).

Therefore, a first empirical estimate is
\[
Q^\pi(x_k,u_k)
\approx
g_k.
\]
If the same state-action pair \((x,u)\) is visited several times in the episode, one averages all
the corresponding returns:
\[
Q^\pi(x,u)
=
\frac{
\sum_{t=0}^{T}\mathbf{1}[x_t=x,\ u_t=u]\,g_t
}{
\sum_{t=0}^{T}\mathbf{1}[x_t=x,\ u_t=u]
}.
\]
This is the Monte Carlo estimate of the Q-function. It says: whenever the episode visits the
state-action pair \((x,u)\), look at the discounted return observed from that point onward; then
average all such observed returns.

This procedure relies on an ergodicity assumption. In order to estimate \(Q^\pi(x,u)\), the
state-action pair \((x,u)\) must be visited sufficiently often. If a pair is never visited, then the
data contain no information about its return.

There are a few important observations.

First, the policy cannot be evaluated until the episode has terminated, because \(g_k\) depends
on all future rewards
\[
r_k,r_{k+1},\dots,r_T.
\]
Thus Monte Carlo policy evaluation is naturally episodic: it computes returns after complete
trajectories have been observed.

Second, the empirical estimate of \(Q^\pi\) does not exactly satisfy the Bellman equation for a
finite amount of data. The Bellman equation is an expectation identity, whereas the Monte Carlo
estimate is based on sample averages. Only in the limit of infinitely many representative samples
do these empirical averages converge to the corresponding expected values.

Third, this method does not exploit the Markov property to make the computation of \(Q^\pi\)
more efficient. Each return \(g_k\) is computed using the full future tail of the episode, rather
than recursively bootstrapping from the value of the next state-action pair. This is one of the
main differences between Monte Carlo learning and temporal-difference methods.

\subsection{Exploration and exploitation}

If we update the policy greedily too early, we may get stuck in a suboptimal behavior. The reason
is simple: at the beginning of learning, the estimate of \(Q^\pi\) is rough. If the policy becomes
greedy with respect to this inaccurate estimate, it may stop trying actions that currently look
bad only because they have not been explored enough.

This is the exploration-exploitation tradeoff. Exploitation means choosing the action that
currently appears best. Exploration means deliberately trying other actions in order to collect
more information.

A common solution is \(\varepsilon\)-greedy policy improvement:
\[
\pi(x,u)
=
\begin{cases}
1-\varepsilon+\dfrac{\varepsilon}{|\cU|}
&
\text{if } u\in\argmax_{\bar u\in\cU}Q(x,\bar u),\\[3mm]
\dfrac{\varepsilon}{|\cU|}
&
\text{otherwise}.
\end{cases}
\]
Equivalently, one can describe the action selection rule as
\[
u
\leftarrow
\begin{cases}
\argmax_{\bar u\in\cU}Q(x,\bar u)
&
\text{with probability }1-\varepsilon,\\
\text{Uniform}(\cU)
&
\text{with probability }\varepsilon.
\end{cases}
\]
A large value of \(\varepsilon\) produces more exploration, while a small value of
\(\varepsilon\) produces more exploitation.

Another common policy improvement rule is the Boltzmann, or softmax, policy:
\[
\pi(x,u)
=
\frac{e^{\beta Q(x,u)}}{\sum_{\bar u\in\cU}e^{\beta Q(x,\bar u)}}.
\]
Here \(\beta\) is an inverse-temperature parameter. If \(\beta\) is small, the probabilities are
close to uniform and the policy explores more. If \(\beta\) is large, the action with the largest
Q-value receives most of the probability mass, and the policy becomes nearly greedy.

As learning progresses from episode to episode, the amount of exploration is usually reduced.
For instance, one may take
\[
\varepsilon\to 0
\qquad\text{or}\qquad
\beta\to\infty.
\]
The intuition is that early episodes should explore the state-action space broadly, whereas later
episodes should exploit the information collected so far. If this is done properly, then with
probability one each relevant state-action pair is visited infinitely often, while the policy
eventually converges to a greedy policy.

\subsection{Scalability and parametrization of Q}

In principle, for a finite MDP, one could store the Q-function in a table with one entry for each
state-action pair. This is the tabular representation:
\[
Q
=
\{Q(x,u):x\in\cX,\ u\in\cU\}.
\]
However, this quickly becomes infeasible when the state and action spaces are large. If the state
has dimension \(n\), the number of possible states typically grows exponentially with \(n\).
Multiple-input systems increase the number of actions, and continuous state or action spaces
make exact tabular storage impossible.

For this reason, the Q-function is very often parametrized by a lower-dimensional vector of
parameters:
\[
Q(x,u)=Q_\theta(x,u),
\qquad
\theta\in\R^d.
\]
The idea is that \(d\) should be much smaller than \(NM\), so that the Q-function can be stored
and learned more efficiently.

This has several advantages:
\begin{itemize}
    \item the memory requirement is reduced from \(NM\) entries to \(d\) parameters;
    \item the problem size is apparently independent of the number of states \(N\), which is
    essential for continuous state spaces;
    \item a good parametrization can generalize from observed state-action pairs to unobserved
    ones, through interpolation or extrapolation;
    \item prior knowledge about the problem can be encoded in the choice of features.
\end{itemize}

There are also important disadvantages:
\begin{itemize}
    \item the number of degrees of freedom is reduced;
    \item the true solution of the Bellman equation for a given policy may not belong to the
    chosen function class
    \[
    \mathcal{Q}
    =
    \{Q_\theta:\theta\in\R^d\};
    \]
    \item convergence of policy iteration and optimality of the limiting policy may be affected;
    \item computing the parameter \(\theta\) that best approximates the data collected from an
    episode may itself be computationally expensive.
\end{itemize}

Thus parametrization is necessary for scalability, but it also introduces approximation error.

\subsection{Linear parametrization}

A simple and important case is linear parametrization. More advanced parametrizations, such as
neural networks, are often used in practice, but the linear case already contains the main idea.

We choose basis functions
\[
\phi_\ell(x,u),
\qquad
\ell=1,\dots,d,
\]
and write
\[
Q_\theta(x,u)
=
\sum_{\ell=1}^{d}\phi_\ell(x,u)\theta_\ell
=
\phi(x,u)^\top\theta,
\]
where
\[
\phi(x,u)
=
\begin{bmatrix}
\phi_1(x,u)\\
\vdots\\
\phi_d(x,u)
\end{bmatrix}.
\]
The functions \(\phi_\ell(x,u)\) may be simple polynomial features such as
\[
x,\quad u,\quad x^2,\quad u^2,\quad xu,\dots,
\]
quadratic forms such as
\[
x^\top A_\ell x,
\]
or problem-specific features. For example, in a game-like environment, one might use quantities
such as total height, number of holes, or other hand-crafted descriptors.

The key point is that the Q-function is no longer represented by one independent value for each
state-action pair. Instead, all Q-values are tied together through the common parameter vector
\(\theta\).

\subsection{Projected Bellman equation}

With parametrization, we cannot generally solve the Bellman equation exactly. For a fixed policy
\(\pi\), the Q-Bellman equation is
\[
Q(x,u)
=
R_x^u
+
\gamma\sum_{x'\in\cX}P^u_{x,x'}
\sum_{u'\in\cU}\pi(x',u')Q(x',u').
\]
Define the Bellman operator
\[
\mathcal{B}(Q)(x,u)
:=
R_x^u
+
\gamma\sum_{x'\in\cX}P^u_{x,x'}
\sum_{u'\in\cU}\pi(x',u')Q(x',u').
\]
The exact Bellman equation is
\[
Q=\mathcal{B}(Q).
\]
However, the function \(\mathcal{B}(Q_\theta)\) will in general not belong to the parametrized
class
\[
\mathcal{Q}
=
\{Q_\theta:\theta\in\R^d\}.
\]
Therefore, after applying the Bellman operator, we project the result back onto the function
class.

Let \(\rho\) be a diagonal \(NM\times NM\) weight matrix. It defines the weighted norm
\[
\norm{Q}_\rho^2
=
Q^\top\rho Q.
\]
The best approximation operator is
\[
\Pi(Q)
=
\argmin_{Q_\theta\in\mathcal{Q}}
\norm{Q_\theta-Q}_\rho^2.
\]
The projected Bellman equation is then
\[
Q_\theta
=
\Pi\bigl(\mathcal{B}(Q_\theta)\bigr).
\]
This equation says that \(Q_\theta\) should be the best approximation, within the chosen
parametrized class, of its Bellman update.

\subsection{Matrix form of the projected Bellman equation}

To write the projected Bellman equation in matrix form, enumerate all state-action pairs and
stack the corresponding values into a vector
\[
Q\in\R^{NM}.
\]
Similarly, define
\[
R\in\R^{NM}
\]
as the vector collecting the expected immediate rewards \(R_x^u\) for all state-action pairs.

Define the state-action transition matrix \(\mathsf{T}^{\pi}\in\R^{NM\times NM}\) by
\[
\mathsf{T}^{\pi}_{(x,u),(x',u')}
=
P^u_{x,x'}\pi(x',u').
\]
This is the transition probability between state-action pairs: starting from \((x,u)\), the system
moves to \(x'\) with probability \(P^u_{x,x'}\), and then policy \(\pi\) selects \(u'\) with
probability \(\pi(x',u')\).

Then the Bellman operator can be written compactly as
\[
\mathcal{B}(Q)
=
R+\gamma\mathsf{T}^{\pi}Q.
\]
For the linear parametrization, collect the feature vectors in the matrix
\[
\Phi
=
\begin{bmatrix}
\phi(x_1,u_1)^\top\\
\phi(x_1,u_2)^\top\\
\vdots\\
\phi(x_N,u_M)^\top
\end{bmatrix}
\in\R^{NM\times d}.
\]
Then
\[
Q_\theta=\Phi\theta.
\]

The projected Bellman equation becomes
\[
\Phi\theta
=
\Pi\left(R+\gamma\mathsf{T}^{\pi}\Phi\theta\right).
\]
The solution is characterized by the normal equations
\[
\Phi^\top\rho\Phi\,\theta
=
\gamma\Phi^\top\rho\mathsf{T}^{\pi}\Phi\,\theta
+
\Phi^\top\rho R.
\]
Equivalently,
\[
\left(
\Phi^\top\rho\Phi
-
\gamma\Phi^\top\rho\mathsf{T}^{\pi}\Phi
\right)\theta
=
\Phi^\top\rho R.
\]
Thus, instead of solving a system with \(NM\) unknown Q-values, we solve a system of \(d\)
linear equations in the parameter vector \(\theta\).

\paragraph{Derivation.}
Since \(Q_\theta=\Phi\theta\), the Bellman update is
\[
\mathcal{B}(Q_\theta)
=
R+\gamma\mathsf{T}^{\pi}\Phi\theta.
\]
The projection step searches for a parameter \(\hat\theta\) minimizing
\[
\norm{\Phi\hat\theta-R-\gamma\mathsf{T}^{\pi}\Phi\theta}_\rho^2.
\]
For a weighted least-squares problem, the normal equations are
\[
\Phi^\top\rho\Phi\,\hat\theta
=
\Phi^\top\rho
\left(
R+\gamma\mathsf{T}^{\pi}\Phi\theta
\right).
\]
At a fixed point of the projected Bellman equation we have \(\hat\theta=\theta\), and therefore
\[
\Phi^\top\rho\Phi\,\theta
=
\gamma\Phi^\top\rho\mathsf{T}^{\pi}\Phi\,\theta
+
\Phi^\top\rho R.
\]
This gives the equation above.

\subsection{Model-free computation of the projected Bellman equation}

The equation
\[
\left(
\Phi^\top\rho\Phi
-
\gamma\Phi^\top\rho\mathsf{T}^{\pi}\Phi
\right)\theta
=
\Phi^\top\rho R
\]
still appears to contain model information through \(\rho\), \(\mathsf{T}^{\pi}\), and \(R\).
The key observation is that these matrices can be estimated from data.

Suppose we observe an episode generated under policy \(\pi\):
\[
(x_0,u_0,r_0),\ (x_1,u_1,r_1),\ \dots,\ (x_T,u_T,r_T).
\]
For each sample, define
\[
\phi_k
:=
\phi(x_k,u_k),
\qquad
\phi_{k+1}
:=
\phi(x_{k+1},u_{k+1}).
\]
Then for every transition in the episode we construct the three quantities
\[
\textcolor{red}{\phi(x_k,u_k)\phi(x_k,u_k)^\top},
\]
\[
\textcolor{blue}{\phi(x_k,u_k)\phi(x_{k+1},u_{k+1})^\top},
\]
and
\[
\textcolor{green!60!black}{\phi(x_k,u_k)r_k}.
\]
Their empirical averages estimate the corresponding matrices:
\[
\widehat M_0
=
\frac{1}{T}\sum_{k=0}^{T-1}
\textcolor{red}{\phi_k\phi_k^\top}
\approx
\Phi^\top\rho\Phi,
\]
\[
\widehat M_1
=
\frac{1}{T}\sum_{k=0}^{T-1}
\textcolor{blue}{\phi_k\phi_{k+1}^\top}
\approx
\Phi^\top\rho\mathsf{T}^{\pi}\Phi,
\]
and
\[
\widehat b
=
\frac{1}{T}\sum_{k=0}^{T-1}
\textcolor{green!60!black}{\phi_k r_k}
\approx
\Phi^\top\rho R.
\]
Therefore, the projected Bellman equation can be approximated from data by
\[
\left(
\widehat M_0
-
\gamma\widehat M_1
\right)\theta
=
\widehat b.
\]
Equivalently,
\[
\left(
\frac{1}{T}\sum_{k=0}^{T-1}\phi_k\phi_k^\top
-
\gamma
\frac{1}{T}\sum_{k=0}^{T-1}\phi_k\phi_{k+1}^\top
\right)\theta
=
\frac{1}{T}\sum_{k=0}^{T-1}\phi_k r_k.
\]

This is a model-free policy evaluation method for linearly parametrized Q-functions. The
transition probabilities are not explicitly known, but their effect is captured through observed
successive pairs
\[
(x_k,u_k)
\quad\text{and}\quad
(x_{k+1},u_{k+1}).
\]
The matrix \(\rho\) is also not chosen explicitly in this empirical version: it is encoded by the
frequency with which the different state-action pairs appear in the dataset.
This is exactly the \((i,j)\)-entry of \(\Phi^\top\rho\Phi\).

The same reasoning applies to the other two terms. The average of
\[
\phi(x_k,u_k)\phi(x_{k+1},u_{k+1})^\top
\]
estimates the transition-weighted matrix
\[
\Phi^\top\rho\mathsf{T}^{\pi}\Phi,
\]
while the average of
\[
\phi(x_k,u_k)r_k
\]
estimates
\[
\Phi^\top\rho R.
\]
Thus, the matrices appearing in the projected Bellman equation can be replaced by empirical
averages computed directly from episodes.

\subsection{Empirical interpretation of the weighted matrices}

Let us make the connection between empirical averages and the weighted matrices even more
explicit. Enumerate the state-action pairs by an index
\[
s=1,\dots,S,
\qquad S=NM.
\]
For each state-action pair \(s\), let
\[
\phi(s)
=
\begin{bmatrix}
\phi_1(s)\\
\vdots\\
\phi_d(s)
\end{bmatrix}
\in\R^d
\]
be the feature vector. The feature matrix is then
\[
\Phi
=
\begin{bmatrix}
\phi(s_1)^\top\\
\phi(s_2)^\top\\
\vdots\\
\phi(s_S)^\top
\end{bmatrix}
\in\R^{S\times d}.
\]
Equivalently, if we denote by \(\phi_j\in\R^S\) the \(j\)-th column of \(\Phi\), then
\[
\Phi
=
\begin{bmatrix}
\phi_1 & \cdots & \phi_d
\end{bmatrix}.
\]

Consider first the matrix
\[
\Phi^\top\rho\Phi.
\]
Since \(\rho\) is diagonal, we can write
\[
\Phi^\top\rho\Phi
=
\begin{bmatrix}
\phi_1^\top\\
\vdots\\
\phi_d^\top
\end{bmatrix}
\begin{bmatrix}
\rho(1) & & \\
& \ddots & \\
& & \rho(S)
\end{bmatrix}
\begin{bmatrix}
\phi_1 & \cdots & \phi_d
\end{bmatrix}.
\]
A generic element in position \((i,j)\) is therefore
\[
\left[\Phi^\top\rho\Phi\right]_{ij}
=
\phi_i^\top
\begin{bmatrix}
\rho(1) & & \\
& \ddots & \\
& & \rho(S)
\end{bmatrix}
\phi_j
=
\sum_{s=1}^{S}\rho(s)\phi_i(s)\phi_j(s).
\]
This is exactly the same expression obtained from the empirical average of the sample matrices
\[
\phi(s)\phi(s)^\top.
\]
Indeed, when many samples are collected, the frequency with which the state-action pair \(s\)
appears in the dataset plays the role of the weight \(\rho(s)\).

A similar reasoning applies to the other two matrices. The empirical average of
\[
\phi(x_k,u_k)\phi(x_{k+1},u_{k+1})^\top
\]
estimates
\[
\Phi^\top\rho\mathsf{T}^{\pi}\Phi,
\]
because it correlates the feature vector of the current state-action pair with the feature vector
of the next state-action pair. Similarly, the empirical average of
\[
\phi(x_k,u_k)r_k
\]
estimates
\[
\Phi^\top\rho R,
\]
because it correlates the current feature vector with the observed immediate reward.

Thus the matrices appearing in the projected Bellman equation can be constructed directly from
an episode, without explicitly estimating either the transition probabilities or the full reward
model.

\subsection{Monte Carlo learning with linear approximation}

We can now combine the previous ideas into a complete learning procedure. The method
alternates between two steps:
\begin{enumerate}
    \item policy evaluation from data;
    \item greedy policy improvement with exploration.
\end{enumerate}

\subsubsection{Policy evaluation}

Fix a policy \(\pi\) and run an experiment, or one simulated episode, under this policy. This
generates a time series
\[
x_k,\ u_k,\ r_k,
\qquad
k=1,\dots,K.
\]
From this episode, we build the empirical estimates
\[
\textcolor{red}{\Phi^\top\rho\Phi}
\approx
\frac{1}{K}\sum_{k=0}^{K-1}
\textcolor{red}{\phi(x_k,u_k)\phi(x_k,u_k)^\top},
\]
\[
\textcolor{blue}{\Phi^\top\rho\mathsf{T}^{\pi}\Phi}
\approx
\frac{1}{K}\sum_{k=0}^{K-1}
\textcolor{blue}{\phi(x_k,u_k)\phi(x_{k+1},u_{k+1})^\top},
\]
and
\[
\textcolor{green!60!black}{\Phi^\top\rho R}
\approx
\frac{1}{K}\sum_{k=0}^{K-1}
\textcolor{green!60!black}{\phi(x_k,u_k)r_k}.
\]
These empirical averages are then inserted into the projected Bellman equation
\[
\left(
\textcolor{red}{\Phi^\top\rho\Phi}
-
\gamma
\textcolor{blue}{\Phi^\top\rho\mathsf{T}^{\pi}\Phi}
\right)\theta
=
\textcolor{green!60!black}{\Phi^\top\rho R}.
\]
Solving this linear system determines the parameter vector \(\theta\), and therefore the
approximate Q-function
\[
Q_\theta(x,u)=\phi(x,u)^\top\theta.
\]

In empirical form, this means solving
\[
\left(
\frac{1}{K}\sum_{k=0}^{K-1}\phi(x_k,u_k)\phi(x_k,u_k)^\top
-
\gamma
\frac{1}{K}\sum_{k=0}^{K-1}\phi(x_k,u_k)\phi(x_{k+1},u_{k+1})^\top
\right)\theta
=
\frac{1}{K}\sum_{k=0}^{K-1}\phi(x_k,u_k)r_k.
\]
The result is an approximate evaluation of the current policy \(\pi\), obtained without knowing
the transition probabilities \(P^u_{x,x'}\).

\subsubsection{Greedy policy improvement}

Once the approximate Q-function has been computed, the policy can be improved greedily. In
the purely greedy case, we would choose
\[
\mu_{\mathrm{new}}(x)
\in
\argmax_{u\in\cU}Q_\theta(x,u).
\]
However, as discussed before, a purely greedy update may stop exploration too early. Therefore,
in learning we typically keep some exploration noise. For example, with an \(\varepsilon\)-greedy
policy we use
\[
u
\leftarrow
\begin{cases}
\displaystyle \argmax_{\bar u\in\cU}Q_\theta(x,\bar u)
&
\text{with probability }1-\varepsilon,\\[2mm]
\text{Uniform}(\cU)
&
\text{with probability }\varepsilon.
\end{cases}
\]
Thus the new policy mostly exploits the current approximation \(Q_\theta\), but still explores
random actions with probability \(\varepsilon\).

The complete scheme is therefore:
\begin{enumerate}
    \item fix the current policy \(\pi\);
    \item collect one or more episodes under \(\pi\);
    \item use the data to estimate the matrices in the projected Bellman equation;
    \item solve the linear system to compute \(\theta\);
    \item define \(Q_\theta(x,u)=\phi(x,u)^\top\theta\);
    \item improve the policy greedily with respect to \(Q_\theta\), while keeping enough
    exploration.
\end{enumerate}

\subsection{Key points}

The Q-function representation is useful because it separates the two parts of policy iteration in
a way that is well suited to learning:
\begin{itemize}
    \item policy improvement becomes trivial, since one only needs to compare the values
    \(Q(x,u)\) for the available actions;
    \item policy evaluation can be delegated to an experiment or simulation.
\end{itemize}

By alternating between policy improvement and episodic experiments, one can progressively
compute better strategies without explicitly identifying the transition probabilities.

However, sufficient exploration must be guaranteed. If the policy becomes too greedy too soon,
some state-action pairs may not be visited often enough, and the learned Q-function may be
poor in precisely the regions where information is missing.

Q-function approximation is essential for scalability. A tabular Q-function requires one value
for every state-action pair, which becomes infeasible in large or continuous problems. A
parametrized approximation
\[
Q_\theta(x,u)
\]
reduces the memory requirement and can generalize across different state-action pairs, but it
also introduces approximation error.

Many variations of this basic idea exist. In the scientific literature, the terminology can vary:
methods of this type are often related to Monte Carlo policy evaluation, generalized policy
iteration, and approximate policy iteration. The common idea is always the same: use data to
evaluate a policy, then use the resulting value or Q-function estimate to improve the policy.